\documentclass[12pt,a4paper]{article}

\usepackage[T2A]{fontenc}
\usepackage[utf8]{inputenc}
\usepackage[russian]{babel}
\usepackage{amsmath,amssymb,amsthm}
\usepackage{mathtools}
\usepackage{enumitem}
\usepackage{listings}
\usepackage{xcolor}
\usepackage{array}
\usepackage{tikz}
\usepackage{forest}
\usepackage{hyperref}

\DeclareMathOperator*{\bigsqcap}{\sqcap}

\hypersetup{
    colorlinks=true,
    linkcolor=blue,
    urlcolor=blue
}

\lstdefinestyle{code}{
    basicstyle=\ttfamily\small,
    keywordstyle=\color{blue},
    commentstyle=\color{gray},
    stringstyle=\color{teal},
    numbers=left,
    numberstyle=\tiny,
    stepnumber=1,
    frame=single,
    breaklines=true,
    tabsize=2
}

\lstset{style=code}

\newtheorem{theorem}{Теорема}
\newtheorem{lemma}{Лемма}
\newtheorem{definition}{Определение}
\newtheorem{example}{Пример}
\newtheorem{remark}{Замечание}

\title{Билет №35 \\ \small Основы построения трансляторов. Структура оптимизирующего транслятора. Промежуточные представления программы. (ИТМО) \\ Основы построения трансляторов. Структура оптимизирующего транслятора. Промежуточные представления программы: последовательность символов, последовательность лексем, синтаксическое дерево, абстрактное синтаксическое дерево. Уровни промежуточного представления: высокий, средний, низкий. Формы промежуточного представления. (СпБГУ)}
\author{Сериков Владислав Александрович}
\date{13 мая 2026}

\begin{document}

\maketitle
\tableofcontents
\newpage

\section{Транслятор и основные понятия}

\begin{definition}
\textbf{Транслятор} --- это программа, преобразующая программу, записанную на одном языке, называемом \textbf{исходным языком}, в эквивалентную программу на другом языке, называемом \textbf{целевым языком}.
\end{definition}

Частными случаями трансляторов являются:
\begin{itemize}
    \item \textbf{компилятор} --- переводит программу целиком, обычно в машинный код, байт-код, ассемблер или промежуточный код;
    \item \textbf{интерпретатор} --- непосредственно выполняет программу, часто совмещая анализ и исполнение;
    \item \textbf{транспилятор} --- переводит программу с одного высокоуровневого языка на другой, например TypeScript $\to$ JavaScript;
    \item \textbf{ассемблер} --- переводит ассемблерный код в машинный код;
    \item \textbf{декомпилятор} --- пытается восстановить высокоуровневое представление из низкоуровневого кода.
\end{itemize}

\begin{definition}
Два фрагмента программы называются \textbf{семантически эквивалентными}, если во всех допустимых контекстах они имеют одинаковое наблюдаемое поведение: вычисляют те же значения, производят те же допустимые побочные эффекты и завершаются или не завершаются одинаковым образом.
\end{definition}

Главная задача транслятора --- построить целевую программу, сохраняющую семантику исходной программы.

\section{Общая структура транслятора}

Классический компилятор удобно делить на три крупные части:
\begin{enumerate}
    \item \textbf{фронтенд} --- анализ исходной программы;
    \item \textbf{мидлэнд} --- промежуточные представления и машинно-независимые оптимизации;
    \item \textbf{бэкенд} --- генерация и машинно-зависимая оптимизация целевого кода.
\end{enumerate}

Общая схема:

\[
\boxed{\text{Исходный текст}}
\to
\boxed{\text{Лексический анализ}}
\to
\boxed{\text{Синтаксический анализ}}
\to
\boxed{\text{Семантический анализ}}
\to
\boxed{\text{IR}}
\to
\boxed{\text{Оптимизации}}
\to
\boxed{\text{Генерация кода}}
\to
\boxed{\text{Целевая программа}}
\]

\subsection{Фронтенд}

Фронтенд зависит от исходного языка. Его основная задача --- проверить корректность программы и построить начальное внутреннее представление.

К основным фазам фронтенда относятся:
\begin{enumerate}
    \item \textbf{лексический анализ};
    \item \textbf{синтаксический анализ};
    \item \textbf{семантический анализ};
    \item построение таблиц символов;
    \item построение начального промежуточного представления.
\end{enumerate}

\subsubsection{Лексический анализ}

\begin{definition}
\textbf{Лексический анализ} --- это фаза трансляции, преобразующая входную последовательность символов в последовательность лексем.
\end{definition}

\begin{definition}
\textbf{Лексема} --- минимальная значимая единица исходной программы: идентификатор, ключевое слово, число, строковый литерал, знак операции, разделитель и т.д.
\end{definition}

Например, исходный фрагмент

\begin{lstlisting}
x = a + 42;
\end{lstlisting}

может быть преобразован в последовательность токенов:

\[
\langle id,x\rangle,\quad
\langle = \rangle,\quad
\langle id,a\rangle,\quad
\langle + \rangle,\quad
\langle number,42\rangle,\quad
\langle ; \rangle.
\]

Лексический анализатор обычно:
\begin{itemize}
    \item удаляет пробелы и комментарии;
    \item распознаёт идентификаторы, литералы и служебные слова;
    \item сообщает о недопустимых символах;
    \item сохраняет позиции токенов для диагностики ошибок.
\end{itemize}

\subsubsection{Синтаксический анализ}

\begin{definition}
\textbf{Синтаксический анализ} --- это фаза, проверяющая, принадлежит ли последовательность лексем грамматике языка, и строящая синтаксическое дерево.
\end{definition}

Обычно синтаксис языка описывается контекстно-свободной грамматикой:

\[
G = (N, \Sigma, P, S),
\]

где:
\begin{itemize}
    \item $N$ --- множество нетерминалов;
    \item $\Sigma$ --- множество терминалов;
    \item $P$ --- множество правил вывода;
    \item $S \in N$ --- стартовый символ.
\end{itemize}

\begin{definition}
\textbf{Синтаксическое дерево}, или \textbf{дерево разбора}, --- дерево, отражающее полный вывод строки в грамматике. Внутренние вершины соответствуют нетерминалам, листья --- терминалам.
\end{definition}

Например, для выражения

\[
a + b * c
\]

и грамматики

\[
E \to E + T \mid T,\qquad
T \to T * F \mid F,\qquad
F \to id
\]

дерево разбора учитывает все нетерминалы $E,T,F$ и показывает приоритет умножения над сложением.

\subsubsection{Семантический анализ}

\begin{definition}
\textbf{Семантический анализ} --- это фаза трансляции, проверяющая свойства программы, которые обычно нельзя выразить чисто синтаксически.
\end{definition}

Типичные задачи семантического анализа:
\begin{itemize}
    \item проверка объявлений идентификаторов;
    \item проверка областей видимости;
    \item проверка типов;
    \item проверка количества и типов аргументов функций;
    \item разрешение перегрузок;
    \item проверка допустимости операторов \texttt{return}, \texttt{break}, \texttt{continue};
    \item построение или уточнение таблицы символов.
\end{itemize}

\begin{definition}
\textbf{Таблица символов} --- структура данных, хранящая информацию об идентификаторах программы: именах переменных, функций, классов, типах, областях видимости, адресах, модификаторах доступа и других свойствах.
\end{definition}

\subsection{Мидлэнд}

Мидлэнд работает с промежуточным представлением программы. Он, как правило, не зависит от исходного языка и в меньшей степени зависит от целевой архитектуры.

Основные задачи мидлэнда:
\begin{itemize}
    \item построение промежуточного кода;
    \item анализ потока управления;
    \item анализ потока данных;
    \item машинно-независимые оптимизации;
    \item преобразование промежуточных представлений.
\end{itemize}

Примеры машинно-независимых оптимизаций:
\begin{itemize}
    \item свёртка констант;
    \item распространение констант;
    \item удаление мёртвого кода;
    \item устранение общих подвыражений;
    \item инвариантное вынесение из циклов;
    \item упрощение алгебраических выражений;
    \item встраивание функций;
    \item преобразование к SSA-форме.
\end{itemize}

\subsection{Бэкенд}

Бэкенд зависит от целевой архитектуры. Его задача --- получить эффективный целевой код.

Основные фазы бэкенда:
\begin{enumerate}
    \item выбор машинных инструкций;
    \item распределение регистров;
    \item планирование инструкций;
    \item генерация ассемблерного или машинного кода;
    \item машинно-зависимые оптимизации.
\end{enumerate}

Примеры машинно-зависимых оптимизаций:
\begin{itemize}
    \item использование специальных инструкций процессора;
    \item оптимизация адресации;
    \item устранение лишних загрузок и сохранений;
    \item переупорядочивание инструкций для конвейера;
    \item оптимизация использования регистров.
\end{itemize}

\section{Структура оптимизирующего транслятора}

Оптимизирующий транслятор не только сохраняет семантику программы, но и пытается улучшить некоторые характеристики целевого кода:
\begin{itemize}
    \item время выполнения;
    \item размер кода;
    \item энергопотребление;
    \item использование памяти;
    \item количество обращений к памяти;
    \item количество ветвлений;
    \item локальность данных.
\end{itemize}

Типовая структура оптимизирующего компилятора:

\[
\begin{array}{c}
\text{Исходная программа}
\\[4pt]
\downarrow
\\[4pt]
\text{Лексический анализ}
\\[4pt]
\downarrow
\\[4pt]
\text{Синтаксический анализ}
\\[4pt]
\downarrow
\\[4pt]
\text{Семантический анализ}
\\[4pt]
\downarrow
\\[4pt]
\text{Высокоуровневое IR}
\\[4pt]
\downarrow
\\[4pt]
\text{Высокоуровневые оптимизации}
\\[4pt]
\downarrow
\\[4pt]
\text{Среднеуровневое IR}
\\[4pt]
\downarrow
\\[4pt]
\text{Машинно-независимые оптимизации}
\\[4pt]
\downarrow
\\[4pt]
\text{Низкоуровневое IR}
\\[4pt]
\downarrow
\\[4pt]
\text{Машинно-зависимые оптимизации}
\\[4pt]
\downarrow
\\[4pt]
\text{Генерация кода}
\\[4pt]
\downarrow
\\[4pt]
\text{Целевая программа}
\end{array}
\]

\subsection{Критерии корректности оптимизаций}

\begin{definition}
\textbf{Оптимизация} --- преобразование программы, сохраняющее её семантику и улучшающее или подготавливающее к улучшению некоторую характеристику программы.
\end{definition}

Не всякая оптимизация обязательно улучшает код в каждом конкретном случае. Например, встраивание функции может уменьшить накладные расходы на вызов, но увеличить размер кода.

Корректная оптимизация должна быть \textbf{семантически сохраняющей}. Например, выражение

\[
x \cdot 2
\]

можно заменить на

\[
x + x
\]

для целочисленных типов без переполнения, но в языках с определённым поведением переполнения или с числами с плавающей точкой такая замена может быть некорректной либо требовать дополнительных условий.

\section{Промежуточные представления программы}

\begin{definition}
\textbf{Промежуточное представление}, или \textbf{IR} от английского \emph{Intermediate Representation}, --- внутренняя форма представления программы, используемая транслятором между исходным текстом и целевым кодом.
\end{definition}

Промежуточные представления нужны для:
\begin{itemize}
    \item разделения фронтенда и бэкенда;
    \item поддержки нескольких исходных языков и нескольких целевых архитектур;
    \item проведения оптимизаций;
    \item упрощения анализа программы;
    \item унификации обработки разных языковых конструкций.
\end{itemize}

Если имеется $m$ исходных языков и $n$ целевых архитектур, то без IR может потребоваться $m \cdot n$ трансляторов. При наличии общего IR достаточно примерно $m+n$ компонентов: $m$ фронтендов и $n$ бэкендов.

\[
\begin{array}{ccc}
L_1 & \searrow & \\
L_2 & \to & IR \to A_1,A_2,\ldots,A_n \\
L_3 & \nearrow &
\end{array}
\]

\section{Последовательные формы представления программы}

\subsection{Последовательность символов}

На самом раннем этапе программа является последовательностью символов:

\[
c_1 c_2 c_3 \ldots c_n.
\]

Например:

\begin{lstlisting}
x = a + 42;
\end{lstlisting}

На этом уровне известны только символы:

\[
\texttt{x},\ \texttt{ },\ \texttt{=},\ \texttt{ },\ \texttt{a},\ \texttt{ },\ \texttt{+},\ \texttt{ },\ \texttt{4},\ \texttt{2},\ \texttt{;}.
\]

Особенности этого уровня:
\begin{itemize}
    \item нет информации о лексемах;
    \item нет информации о синтаксической структуре;
    \item пробелы и комментарии ещё присутствуют;
    \item ошибки связаны с недопустимыми символами или некорректной кодировкой.
\end{itemize}

\subsection{Последовательность лексем}

После лексического анализа программа превращается в последовательность токенов:

\[
t_1,t_2,\ldots,t_k.
\]

Для фрагмента

\begin{lstlisting}
x = a + 42;
\end{lstlisting}

можно получить:

\[
\langle ID,x\rangle,\
\langle ASSIGN\rangle,\
\langle ID,a\rangle,\
\langle PLUS\rangle,\
\langle NUM,42\rangle,\
\langle SEMICOLON\rangle.
\]

На этом уровне уже понятны:
\begin{itemize}
    \item границы идентификаторов;
    \item числовые и строковые литералы;
    \item ключевые слова;
    \item операторы;
    \item разделители.
\end{itemize}

Но ещё не построены:
\begin{itemize}
    \item дерево выражений;
    \item структура операторов;
    \item области видимости;
    \item типы выражений.
\end{itemize}

\subsection{Синтаксическое дерево}

\begin{definition}
\textbf{Синтаксическое дерево} --- дерево, отражающее структуру вывода программы в грамматике.
\end{definition}

Для выражения

\[
a + b * c
\]

полное синтаксическое дерево может иметь вид:

\[
\begin{forest}
[E
    [E
        [T
            [F
                [id [a]]
            ]
        ]
    ]
    [+]
    [T
        [T
            [F
                [id [b]]
            ]
        ]
        [*]
        [F
            [id [c]]
        ]
    ]
]
\end{forest}
\]

Такое дерево содержит много служебных нетерминалов, нужных для описания грамматики, но не всегда важных для дальнейших фаз компилятора.

\subsection{Абстрактное синтаксическое дерево}

\begin{definition}
\textbf{Абстрактное синтаксическое дерево}, или \textbf{AST}, --- дерево, отражающее существенную синтаксическую структуру программы без лишних деталей конкретной грамматики.
\end{definition}

Для выражения

\[
a + b * c
\]

AST имеет вид:

\[
\begin{forest}
[+
    [a]
    [*
        [b]
        [c]
    ]
]
\end{forest}
\]

В AST обычно отсутствуют:
\begin{itemize}
    \item лишние нетерминалы;
    \item скобки, если их роль уже отражена структурой дерева;
    \item точки с запятой;
    \item некоторые служебные токены;
    \item детали конкретного синтаксиса, не влияющие на семантику.
\end{itemize}

AST удобно использовать для:
\begin{itemize}
    \item семантического анализа;
    \item проверки типов;
    \item преобразования исходных конструкций;
    \item генерации более низкоуровневого IR.
\end{itemize}

\section{Уровни промежуточного представления}

Промежуточные представления часто делят на:
\begin{enumerate}
    \item высокоуровневые;
    \item среднеуровневые;
    \item низкоуровневые.
\end{enumerate}

\subsection{Высокоуровневое промежуточное представление}

\begin{definition}
\textbf{Высокоуровневое IR} --- представление, близкое к исходному языку и сохраняющее многие его конструкции.
\end{definition}

Примеры высокоуровневых конструкций:
\begin{itemize}
    \item циклы \texttt{for}, \texttt{while};
    \item условные операторы;
    \item вызовы функций;
    \item исключения;
    \item объекты и методы;
    \item замыкания;
    \item шаблоны или обобщённые типы;
    \item массивы и срезы.
\end{itemize}

Пример исходной программы:

\begin{lstlisting}
for (int i = 0; i < n; i++) {
    s += a[i];
}
\end{lstlisting}

Высокоуровневое IR может сохранять цикл как специальную конструкцию:

\begin{lstlisting}
For(
    init: i = 0,
    cond: i < n,
    step: i = i + 1,
    body: s = s + a[i]
)
\end{lstlisting}

Преимущества:
\begin{itemize}
    \item удобно анализировать исходные конструкции;
    \item можно проводить оптимизации высокого уровня;
    \item сохраняется информация о типах и областях видимости.
\end{itemize}

Недостатки:
\begin{itemize}
    \item сложно напрямую генерировать машинный код;
    \item представление зависит от особенностей исходного языка;
    \item сложнее унифицировать несколько языков.
\end{itemize}

\subsection{Среднеуровневое промежуточное представление}

\begin{definition}
\textbf{Среднеуровневое IR} --- представление, достаточно абстрактное для машинно-независимых оптимизаций, но уже не содержащее многих сложных конструкций исходного языка.
\end{definition}

На этом уровне обычно:
\begin{itemize}
    \item циклы выражаются через переходы и базовые блоки;
    \item сложные выражения разбиваются на простые операции;
    \item появляются временные переменные;
    \item явно задаётся поток управления;
    \item может использоваться SSA-форма.
\end{itemize}

Например:

\begin{lstlisting}
s = 0;
i = 0;
while (i < n) {
    s = s + a[i];
    i = i + 1;
}
\end{lstlisting}

может быть представлено как:

\begin{lstlisting}
B1:
    s = 0
    i = 0
    goto B2

B2:
    t1 = i < n
    if t1 goto B3 else B4

B3:
    t2 = a[i]
    s = s + t2
    i = i + 1
    goto B2

B4:
    return s
\end{lstlisting}

Преимущества:
\begin{itemize}
    \item удобно строить граф потока управления;
    \item удобно проводить анализ потока данных;
    \item подходит для многих машинно-независимых оптимизаций;
    \item сравнительно не зависит от исходного языка.
\end{itemize}

\subsection{Низкоуровневое промежуточное представление}

\begin{definition}
\textbf{Низкоуровневое IR} --- представление, близкое к машинным инструкциям, но ещё сохраняющее некоторую абстракцию от конкретного процессора.
\end{definition}

На низком уровне обычно явно представлены:
\begin{itemize}
    \item операции загрузки и сохранения;
    \item адреса;
    \item регистры или виртуальные регистры;
    \item машинные типы;
    \item вызовы функций по соглашению вызова;
    \item переходы и метки.
\end{itemize}

Пример:

\begin{lstlisting}
load r1, [a + i * 4]
add  r2, s, r1
mov  s, r2
add  i, i, 1
cmp  i, n
jl   B3
\end{lstlisting}

Преимущества:
\begin{itemize}
    \item близко к генерации машинного кода;
    \item удобно выполнять машинно-зависимые оптимизации;
    \item удобно проводить распределение регистров.
\end{itemize}

Недостатки:
\begin{itemize}
    \item теряется высокоуровневая информация;
    \item сложнее проводить оптимизации, зависящие от структуры исходной программы.
\end{itemize}

\section{Формы промежуточного представления}

Промежуточные представления можно классифицировать не только по уровню абстракции, но и по форме.

\subsection{Деревья}

Деревья хорошо отражают вложенную структуру выражений.

Пример выражения:

\[
x = (a+b)\cdot c
\]

AST:

\[
\begin{forest}
[=
    [x]
    [*
        [+
            [a]
            [b]
        ]
        [c]
    ]
]
\end{forest}
\]

Деревья удобны для:
\begin{itemize}
    \item синтаксического анализа;
    \item семантических проверок;
    \item локальных преобразований выражений.
\end{itemize}

Но деревья хуже подходят для:
\begin{itemize}
    \item анализа циклов;
    \item анализа потока управления;
    \item глобальных оптимизаций;
    \item генерации низкоуровневого кода.
\end{itemize}

\subsection{Трёхадресный код}

\begin{definition}
\textbf{Трёхадресный код} --- линейное промежуточное представление, в котором каждая инструкция содержит не более одного оператора и обычно имеет вид:
\[
x = y \operatorname{op} z.
\]
\end{definition}

Пример:

\[
x = (a+b)\cdot(c-d)
\]

переводится в:

\begin{lstlisting}
t1 = a + b
t2 = c - d
t3 = t1 * t2
x  = t3
\end{lstlisting}

Основные виды инструкций трёхадресного кода:
\begin{itemize}
    \item \texttt{x = y op z};
    \item \texttt{x = op y};
    \item \texttt{x = y};
    \item \texttt{goto L};
    \item \texttt{if x relop y goto L};
    \item \texttt{param x};
    \item \texttt{call f, n};
    \item \texttt{return x};
    \item \texttt{x = y[i]};
    \item \texttt{x[i] = y};
    \item \texttt{x = *p};
    \item \texttt{*p = x}.
\end{itemize}

\subsection{Квадрупли, трипли и косвенные трипли}

Трёхадресный код может храниться в разных структурах.

\subsubsection{Квадрупли}

\begin{definition}
\textbf{Квадрупль} --- запись вида
\[
(op, arg_1, arg_2, result).
\]
\end{definition}

Для кода

\begin{lstlisting}
t1 = a + b
t2 = c - d
t3 = t1 * t2
x  = t3
\end{lstlisting}

получим:

\[
\begin{array}{c|c|c|c|c}
№ & op & arg_1 & arg_2 & result \\
\hline
1 & + & a & b & t1 \\
2 & - & c & d & t2 \\
3 & * & t1 & t2 & t3 \\
4 & = & t3 & - & x
\end{array}
\]

\subsubsection{Трипли}

\begin{definition}
\textbf{Трипль} --- запись вида
\[
(op, arg_1, arg_2),
\]
где результат операции не имеет отдельного имени, а обозначается номером самой инструкции.
\end{definition}

\[
\begin{array}{c|c|c|c}
№ & op & arg_1 & arg_2 \\
\hline
1 & + & a & b \\
2 & - & c & d \\
3 & * & (1) & (2) \\
4 & = & (3) & x
\end{array}
\]

\subsubsection{Косвенные трипли}

\begin{definition}
\textbf{Косвенные трипли} используют дополнительную таблицу указателей на трипли, что упрощает перестановку инструкций.
\end{definition}

Например:

\[
\begin{array}{c|c}
Позиция & Указатель \\
\hline
1 & \#1 \\
2 & \#2 \\
3 & \#3 \\
4 & \#4
\end{array}
\]

При оптимизациях можно менять порядок указателей, не перемещая сами трипли.

\subsection{Стековое представление}

В стековом IR операнды передаются через стек.

Пример:

\[
x = (a+b)c
\]

Стековый код:

\begin{lstlisting}
push a
push b
add
push c
mul
pop x
\end{lstlisting}

Такое представление используется в виртуальных машинах, например в JVM и некоторых байт-кодах.

\subsection{Графовые представления}

К важным графовым представлениям относятся:
\begin{itemize}
    \item граф потока управления;
    \item граф потока данных;
    \item граф зависимостей;
    \item граф вызовов;
    \item SSA-граф.
\end{itemize}

\section{Граф потока управления}

\begin{definition}
\textbf{Базовый блок} --- максимальная последовательность инструкций, в которую управление входит только через первую инструкцию, а выходит только из последней.
\end{definition}

\begin{definition}
\textbf{Граф потока управления}, или \textbf{CFG} от английского \emph{Control Flow Graph}, --- ориентированный граф, вершинами которого являются базовые блоки, а рёбрами --- возможные передачи управления между блоками.
\end{definition}

Пример:

\begin{lstlisting}
if (x > 0)
    y = 1;
else
    y = 2;
z = y + 3;
\end{lstlisting}

Трёхадресный код:

\begin{lstlisting}
B1:
    if x > 0 goto B2 else B3

B2:
    y = 1
    goto B4

B3:
    y = 2
    goto B4

B4:
    z = y + 3
\end{lstlisting}

CFG:

\[
B1 \to B2,\qquad B1 \to B3,\qquad B2 \to B4,\qquad B3 \to B4.
\]

\section{Алгоритм построения базовых блоков}

\subsection{Идея}

Чтобы построить CFG по линейному трёхадресному коду, сначала выделяют лидеры базовых блоков.

\begin{definition}
\textbf{Лидер} --- первая инструкция некоторого базового блока.
\end{definition}

\subsection{Алгоритм}

\textbf{Вход:} последовательность инструкций трёхадресного кода.

\textbf{Выход:} множество базовых блоков.

\begin{enumerate}
    \item Первая инструкция программы является лидером.
    \item Любая инструкция, на которую есть переход, является лидером.
    \item Любая инструкция, непосредственно следующая за условным или безусловным переходом, является лидером.
    \item Каждый базовый блок начинается с лидера и продолжается до инструкции перед следующим лидером или до конца программы.
\end{enumerate}

\subsection{Минимальный пример}

Пусть дан код:

\begin{lstlisting}
1: a = 1
2: if x < 0 goto 5
3: a = 2
4: goto 6
5: a = 3
6: b = a + 1
\end{lstlisting}

Шаги алгоритма:
\begin{enumerate}
    \item Инструкция 1 --- лидер, потому что она первая.
    \item Инструкция 5 --- лидер, потому что на неё есть переход из инструкции 2.
    \item Инструкция 6 --- лидер, потому что на неё есть переход из инструкции 4.
    \item Инструкция 3 --- лидер, потому что она следует сразу после условного перехода в инструкции 2.
\end{enumerate}

Множество лидеров:

\[
\{1,3,5,6\}.
\]

Базовые блоки:

\begin{lstlisting}
B1:
1: a = 1
2: if x < 0 goto 5

B2:
3: a = 2
4: goto 6

B3:
5: a = 3

B4:
6: b = a + 1
\end{lstlisting}

Рёбра CFG:

\[
B1 \to B2,\quad B1 \to B3,\quad B2 \to B4,\quad B3 \to B4.
\]

\section{SSA-форма}

\begin{definition}
\textbf{SSA-форма}, или \textbf{Static Single Assignment}, --- промежуточное представление, в котором каждая переменная присваивается ровно один раз.
\end{definition}

Если переменная получает значения из разных ветвей управления, используется специальная $\varphi$-функция.

Пример исходного кода:

\begin{lstlisting}
if (c)
    x = 1;
else
    x = 2;
y = x + 3;
\end{lstlisting}

SSA-форма:

\begin{lstlisting}
if c goto B1 else B2

B1:
    x1 = 1
    goto B3

B2:
    x2 = 2
    goto B3

B3:
    x3 = phi(x1, x2)
    y1 = x3 + 3
\end{lstlisting}

\begin{definition}
$\varphi$-функция выбирает значение в зависимости от того, из какого предшествующего базового блока пришло управление.
\end{definition}

В приведённом примере:
\begin{itemize}
    \item если управление пришло из блока $B1$, то $x_3 = x_1$;
    \item если управление пришло из блока $B2$, то $x_3 = x_2$.
\end{itemize}

\section{Оптимизации промежуточного кода}

\subsection{Свёртка констант}

\begin{definition}
\textbf{Свёртка констант} --- оптимизация, заменяющая выражения, состоящие из констант, на результат их вычисления во время компиляции.
\end{definition}

Пример:

\begin{lstlisting}
x = 2 * 3 + 4
\end{lstlisting}

после оптимизации:

\begin{lstlisting}
x = 10
\end{lstlisting}

\subsubsection{Алгоритм свёртки констант}

\textbf{Вход:} AST или трёхадресный код.

\textbf{Выход:} программа с вычисленными константными выражениями.

\begin{enumerate}
    \item Обойти выражение снизу вверх.
    \item Для каждого узла операции проверить, являются ли все его операнды константами.
    \item Если да, вычислить результат операции.
    \item Заменить узел операции на константный узел.
    \item Повторять до неподвижной точки, если оптимизация выполняется совместно с другими преобразованиями.
\end{enumerate}

Пример на дереве:

\[
\begin{forest}
[+
    [*
        [2]
        [3]
    ]
    [4]
]
\end{forest}
\]

Шаг 1:

\[
2 * 3 \to 6.
\]

Получаем:

\[
\begin{forest}
[+
    [6]
    [4]
]
\end{forest}
\]

Шаг 2:

\[
6 + 4 \to 10.
\]

Итог:

\[
\begin{forest}
[10]
\end{forest}
\]

\subsection{Распространение констант}

\begin{definition}
\textbf{Распространение констант} --- оптимизация, заменяющая использование переменной на константу, если известно, что в данной точке программы переменная имеет это постоянное значение.
\end{definition}

Пример:

\begin{lstlisting}
x = 5
y = x + 1
\end{lstlisting}

после распространения константы:

\begin{lstlisting}
x = 5
y = 5 + 1
\end{lstlisting}

после свёртки констант:

\begin{lstlisting}
x = 5
y = 6
\end{lstlisting}

\subsection{Удаление мёртвого кода}

\begin{definition}
\textbf{Мёртвый код} --- код, результат которого не влияет на наблюдаемое поведение программы.
\end{definition}

Пример:

\begin{lstlisting}
x = 1
x = 2
print(x)
\end{lstlisting}

Присваивание \texttt{x = 1} является мёртвым, если между ним и \texttt{x = 2} значение \texttt{x} не используется.

После удаления:

\begin{lstlisting}
x = 2
print(x)
\end{lstlisting}

\subsection{Устранение общих подвыражений}

\begin{definition}
\textbf{Общее подвыражение} --- выражение, которое уже было вычислено ранее и операнды которого с тех пор не изменились.
\end{definition}

Пример:

\begin{lstlisting}
x = a * b + c
y = a * b + d
\end{lstlisting}

После оптимизации:

\begin{lstlisting}
t = a * b
x = t + c
y = t + d
\end{lstlisting}

\subsection{Инвариантное вынесение из цикла}

\begin{definition}
\textbf{Инвариант цикла} --- выражение, значение которого не меняется от итерации к итерации цикла.
\end{definition}

Пример:

\begin{lstlisting}
while (i < n) {
    x = a * b;
    s = s + x + i;
    i = i + 1;
}
\end{lstlisting}

Если $a$ и $b$ не меняются в цикле, можно преобразовать:

\begin{lstlisting}
x = a * b;
while (i < n) {
    s = s + x + i;
    i = i + 1;
}
\end{lstlisting}

\section{Анализ потока данных}

Многие оптимизации требуют информации о том, какие определения переменных достигают данной точки, какие переменные живы, какие выражения доступны и т.д.

\begin{definition}
\textbf{Анализ потока данных} --- метод статического анализа, при котором для каждой точки программы вычисляется информация о возможных значениях или свойствах программы.
\end{definition}

Типичные задачи:
\begin{itemize}
    \item достигающие определения;
    \item живые переменные;
    \item доступные выражения;
    \item распространение констант;
    \item анализ указателей;
    \item анализ алиасов.
\end{itemize}

\section{Живые переменные}

\begin{definition}
Переменная $x$ называется \textbf{живой} в точке программы, если существует путь выполнения от этой точки до использования значения $x$, на котором $x$ не переопределяется.
\end{definition}

Анализ живых переменных является обратным анализом потока данных.

Для базового блока $B$ определяются множества:
\begin{itemize}
    \item $use[B]$ --- переменные, используемые в $B$ до первого определения в $B$;
    \item $def[B]$ --- переменные, определяемые в $B$;
    \item $in[B]$ --- переменные, живые на входе в $B$;
    \item $out[B]$ --- переменные, живые на выходе из $B$.
\end{itemize}

Уравнения:

\[
out[B] = \bigcup_{S \in succ[B]} in[S],
\]

\[
in[B] = use[B] \cup (out[B] \setminus def[B]).
\]

\subsection{Алгоритм анализа живых переменных}

\textbf{Вход:} CFG программы.

\textbf{Выход:} множества $in[B]$ и $out[B]$ для каждого блока.

\begin{enumerate}
    \item Для каждого блока вычислить $use[B]$ и $def[B]$.
    \item Инициализировать:
    \[
    in[B] = \varnothing,\qquad out[B] = \varnothing.
    \]
    \item Повторять до неподвижной точки:
    \begin{enumerate}
        \item для каждого блока $B$ вычислить:
        \[
        out[B] := \bigcup_{S \in succ[B]} in[S];
        \]
        \item затем вычислить:
        \[
        in[B] := use[B] \cup (out[B] \setminus def[B]).
        \]
    \end{enumerate}
\end{enumerate}

\subsection{Минимальный пример}

Пусть есть блоки:

\begin{lstlisting}
B1:
    a = 1
    b = 2

B2:
    c = a + b

B3:
    print(c)
\end{lstlisting}

Рёбра:

\[
B1 \to B2,\qquad B2 \to B3.
\]

Множества:

\[
use[B1]=\varnothing,\quad def[B1]=\{a,b\},
\]

\[
use[B2]=\{a,b\},\quad def[B2]=\{c\},
\]

\[
use[B3]=\{c\},\quad def[B3]=\varnothing.
\]

Итерации:

\[
in[B3]=\{c\},\quad out[B3]=\varnothing.
\]

\[
out[B2]=in[B3]=\{c\},
\]

\[
in[B2]=\{a,b\}\cup(\{c\}\setminus\{c\})=\{a,b\}.
\]

\[
out[B1]=in[B2]=\{a,b\},
\]

\[
in[B1]=\varnothing\cup(\{a,b\}\setminus\{a,b\})=\varnothing.
\]

Итог:

\[
\begin{array}{c|c|c}
B & in[B] & out[B] \\
\hline
B1 & \varnothing & \{a,b\} \\
B2 & \{a,b\} & \{c\} \\
B3 & \{c\} & \varnothing
\end{array}
\]

\section{Теоретическая основа итерационных анализов}

Для анализа потока данных часто используется теория решёток и неподвижных точек.

\begin{definition}
Частично упорядоченное множество $(L,\sqsubseteq)$ называется \textbf{полной решёткой}, если любое подмножество $S \subseteq L$ имеет точную верхнюю грань $\bigsqcup S$ и точную нижнюю грань $\bigsqcap S$.
\end{definition}

\begin{definition}
Функция $f : L \to L$ называется \textbf{монотонной}, если
\[
x \sqsubseteq y \Rightarrow f(x) \sqsubseteq f(y).
\]
\end{definition}

\begin{definition}
Элемент $x \in L$ называется \textbf{неподвижной точкой} функции $f : L \to L$, если
\[
f(x)=x.
\]
\end{definition}

\begin{theorem}[Теорема Кнастера--Тарского]
Пусть $(L,\sqsubseteq)$ --- полная решётка, а $f:L\to L$ --- монотонная функция. Тогда множество неподвижных точек функции $f$ непусто и само образует полную решётку. В частности, у $f$ существуют наименьшая и наибольшая неподвижные точки:
\[
\operatorname{lfp}(f)=\bigsqcap \{x\in L \mid f(x)\sqsubseteq x\},
\]
\[
\operatorname{gfp}(f)=\bigsqcup \{x\in L \mid x\sqsubseteq f(x)\}.
\]
\end{theorem}

\begin{proof}
Рассмотрим множество префиксных точек:
\[
P=\{x\in L \mid f(x)\sqsubseteq x\}.
\]
Так как $L$ --- полная решётка, множество $P$ имеет точную нижнюю грань:
\[
p=\bigsqcap P.
\]
Покажем, что $p$ --- неподвижная точка.

Так как $p \sqsubseteq x$ для любого $x\in P$, по монотонности $f$ имеем:
\[
f(p)\sqsubseteq f(x).
\]
Но для $x\in P$ выполнено $f(x)\sqsubseteq x$, следовательно:
\[
f(p)\sqsubseteq x
\]
для всех $x\in P$. Значит, $f(p)$ является нижней гранью множества $P$, а поскольку $p$ --- точная нижняя грань, получаем:
\[
f(p)\sqsubseteq p.
\]
Следовательно, $p\in P$.

Теперь применим монотонность к неравенству $f(p)\sqsubseteq p$:
\[
f(f(p))\sqsubseteq f(p).
\]
Значит, $f(p)\in P$. Но $p$ --- нижняя грань $P$, поэтому:
\[
p\sqsubseteq f(p).
\]
Совмещая два неравенства:
\[
f(p)\sqsubseteq p
\quad \text{и} \quad
p\sqsubseteq f(p),
\]
получаем:
\[
f(p)=p.
\]
Следовательно, $p$ --- неподвижная точка.

Более того, для любой неподвижной точки $y$ имеем $f(y)=y$, следовательно $y\in P$. Поскольку $p=\bigsqcap P$, то
\[
p\sqsubseteq y.
\]
Значит, $p$ является наименьшей неподвижной точкой.

Аналогично рассматривается множество постфиксных точек:
\[
Q=\{x\in L \mid x\sqsubseteq f(x)\}.
\]
Его точная верхняя грань
\[
q=\bigsqcup Q
\]
оказывается наибольшей неподвижной точкой. Множество всех неподвижных точек также замкнуто относительно взятия точных верхних и нижних граней внутри соответствующего множества, поэтому оно образует полную решётку.
\end{proof}

\begin{remark}
В компиляторах эта теорема объясняет, почему многие задачи анализа потока данных можно решать как поиск неподвижной точки системы монотонных уравнений на конечной решётке.
\end{remark}

\begin{theorem}[Завершимость итерационного анализа на конечной решётке]
Пусть $L$ --- конечная решётка, а $f:L\to L$ --- монотонная функция. Тогда итерационный процесс
\[
x_0=\bot,\qquad x_{i+1}=f(x_i)
\]
при условии, что $x_i\sqsubseteq x_{i+1}$ для всех $i$, завершается за конечное число шагов и достигает неподвижной точки.
\end{theorem}

\begin{proof}
Так как
\[
x_0\sqsubseteq x_1\sqsubseteq x_2\sqsubseteq \ldots,
\]
получается возрастающая цепь элементов решётки $L$.

Поскольку $L$ конечна, в ней не может существовать бесконечной строго возрастающей цепи. Следовательно, найдётся номер $k$, такой что:
\[
x_k=x_{k+1}.
\]
Но по определению итераций:
\[
x_{k+1}=f(x_k).
\]
Значит:
\[
x_k=f(x_k).
\]
Следовательно, $x_k$ является неподвижной точкой функции $f$, и процесс завершается за конечное число шагов.
\end{proof}

\section{Синтаксический анализ: основные алгоритмические идеи}

Синтаксические анализаторы делятся на:
\begin{itemize}
    \item нисходящие;
    \item восходящие.
\end{itemize}

\subsection{Рекурсивный спуск}

\begin{definition}
\textbf{Рекурсивный спуск} --- метод нисходящего синтаксического анализа, при котором каждому нетерминалу грамматики соответствует процедура.
\end{definition}

Например, для грамматики:

\[
E \to T E',
\]

\[
E' \to + T E' \mid \varepsilon,
\]

\[
T \to id
\]

можно написать процедуры:
\begin{itemize}
    \item \texttt{parseE};
    \item \texttt{parseEPrime};
    \item \texttt{parseT}.
\end{itemize}

\subsubsection{Алгоритм рекурсивного спуска}

\begin{enumerate}
    \item Начать с процедуры стартового нетерминала.
    \item Для текущего нетерминала выбрать правило по текущей входной лексеме.
    \item Если встречается терминал, сравнить его с текущей лексемой.
    \item Если терминал совпал, перейти к следующей лексеме.
    \item Если встречается нетерминал, вызвать соответствующую процедуру.
    \item Если ни одно правило не подходит, сообщить о синтаксической ошибке.
    \item Если стартовая процедура завершилась и вход полностью прочитан, программа синтаксически корректна.
\end{enumerate}

\subsubsection{Минимальный пример}

Грамматика:

\[
E \to T E',
\]

\[
E' \to + T E' \mid \varepsilon,
\]

\[
T \to id.
\]

Вход:

\[
id + id.
\]

Шаги:
\begin{enumerate}
    \item Вызов $E$.
    \item По правилу $E \to T E'$ вызывается $T$.
    \item $T$ распознаёт первый $id$.
    \item Вызывается $E'$.
    \item Текущий токен $+$, выбирается правило $E'\to +TE'$.
    \item Распознаётся $+$.
    \item $T$ распознаёт второй $id$.
    \item Снова вызывается $E'$.
    \item Вход закончился, выбирается $E'\to\varepsilon$.
    \item Разбор успешен.
\end{enumerate}

\section{Лексический анализ и регулярные языки}

Лексемы большинства языков программирования задаются регулярными выражениями.

Например:
\[
ID = [a-zA-Z\_][a-zA-Z0-9\_]^{*},
\]
\[
NUM = [0-9]+.
\]

\begin{definition}
\textbf{Детерминированный конечный автомат} --- пятёрка
\[
A=(Q,\Sigma,\delta,q_0,F),
\]
где:
\begin{itemize}
    \item $Q$ --- конечное множество состояний;
    \item $\Sigma$ --- входной алфавит;
    \item $\delta:Q\times\Sigma\to Q$ --- функция переходов;
    \item $q_0\in Q$ --- начальное состояние;
    \item $F\subseteq Q$ --- множество допускающих состояний.
\end{itemize}
\end{definition}

Лексер можно рассматривать как набор конечных автоматов, распознающих разные классы лексем.

\section{Связь регулярных выражений и конечных автоматов}

\begin{theorem}[Клини]
Язык является регулярным тогда и только тогда, когда он распознаётся конечным автоматом.
\end{theorem}

\begin{proof}
Докажем в двух направлениях.

\textbf{1. Всякий язык, задаваемый регулярным выражением, распознаётся конечным автоматом.}

Доказательство проводится структурной индукцией по построению регулярного выражения.

Базовые случаи:
\begin{itemize}
    \item выражение $\varnothing$ задаёт пустой язык; его распознаёт автомат без достижимых допускающих состояний;
    \item выражение $\varepsilon$ задаёт язык $\{\varepsilon\}$; его распознаёт автомат, у которого начальное состояние является допускающим;
    \item выражение $a$, где $a\in\Sigma$, задаёт язык $\{a\}$; его распознаёт автомат с переходом из начального состояния в допускающее по символу $a$.
\end{itemize}

Индукционный шаг. Пусть для регулярных выражений $r$ и $s$ уже построены конечные автоматы, распознающие языки $L(r)$ и $L(s)$.

Для выражения $r|s$ строится автомат, который из нового начального состояния $\varepsilon$-переходами переходит в начальные состояния автоматов для $r$ и $s$, а допускающими делает их допускающие состояния. Он распознаёт объединение:
\[
L(r)\cup L(s).
\]

Для выражения $rs$ соединяем допускающие состояния автомата для $r$ с начальным состоянием автомата для $s$ $\varepsilon$-переходами. Получаем язык конкатенации:
\[
L(r)L(s).
\]

Для выражения $r^*$ добавляем новое начальное допускающее состояние, $\varepsilon$-переход к автомату для $r$ и $\varepsilon$-переходы из допускающих состояний автомата для $r$ обратно к его начальному состоянию. Получаем язык:
\[
L(r)^*.
\]

Следовательно, для любого регулярного выражения можно построить недетерминированный конечный автомат. Так как всякий недетерминированный конечный автомат можно детерминировать методом подмножеств, язык распознаётся детерминированным конечным автоматом.

\textbf{2. Всякий язык, распознаваемый конечным автоматом, задаётся регулярным выражением.}

Пусть дан конечный автомат. Построим для него обобщённый конечный автомат, в котором рёбра помечены регулярными выражениями. Добавим новое начальное состояние и новое единственное допускающее состояние. Далее будем последовательно удалять промежуточные состояния.

Пусть удаляется состояние $k$. Для каждой пары состояний $i,j$ обновим метку ребра $i\to j$ по формуле:
\[
R_{ij} := R_{ij} \mid R_{ik}(R_{kk})^*R_{kj}.
\]
Эта формула учитывает все пути из $i$ в $j$, которые проходят через $k$ любое число раз.

После удаления всех промежуточных состояний останется одно ребро из нового начального состояния в новое допускающее состояние. Метка этого ребра является регулярным выражением, задающим ровно язык исходного автомата.

Следовательно, каждый язык, распознаваемый конечным автоматом, является регулярным.

Из двух направлений следует утверждение теоремы.
\end{proof}

\section{Понижение представлений}

\begin{definition}
\textbf{Понижение}, или \textbf{lowering}, --- преобразование более высокоуровневого промежуточного представления в более низкоуровневое.
\end{definition}

Пример: цикл \texttt{while}

\begin{lstlisting}
while (x < 10) {
    x = x + 1;
}
\end{lstlisting}

можно понизить до переходов:

\begin{lstlisting}
L1:
    t1 = x < 10
    if t1 goto L2 else L3

L2:
    x = x + 1
    goto L1

L3:
\end{lstlisting}

Понижение позволяет постепенно избавляться от сложных конструкций исходного языка.

\section{Пример полного пути трансляции}

Рассмотрим фрагмент:

\begin{lstlisting}
x = (a + b) * 2;
\end{lstlisting}

\subsection{Последовательность символов}

\[
\texttt{x = (a + b) * 2;}
\]

\subsection{Последовательность лексем}

\[
\langle ID,x\rangle,\
\langle ASSIGN\rangle,\
\langle LPAREN\rangle,\
\langle ID,a\rangle,\
\langle PLUS\rangle,\
\langle ID,b\rangle,\
\langle RPAREN\rangle,\
\langle MUL\rangle,\
\langle NUM,2\rangle,\
\langle SEMICOLON\rangle.
\]

\subsection{AST}

\[
\begin{forest}
[=
    [x]
    [*
        [+
            [a]
            [b]
        ]
        [2]
    ]
]
\end{forest}
\]

\subsection{Трёхадресный код}

\begin{lstlisting}
t1 = a + b
t2 = t1 * 2
x  = t2
\end{lstlisting}

\subsection{Возможная оптимизация}

Если целевая архитектура эффективно выполняет сдвиг, можно заменить умножение на 2:

\begin{lstlisting}
t1 = a + b
t2 = t1 << 1
x  = t2
\end{lstlisting}

Но такая оптимизация требует учёта типа данных и правил переполнения.

\subsection{Низкоуровневое представление}

\begin{lstlisting}
load r1, [a]
load r2, [b]
add  r3, r1, r2
shl  r4, r3, 1
store [x], r4
\end{lstlisting}

\section{Ошибки в трансляторах}

Ошибки, обнаруживаемые транслятором, можно разделить на несколько групп.

\subsection{Лексические ошибки}

Пример:

\begin{lstlisting}
x = 12@3;
\end{lstlisting}

Символ \texttt{@} может быть недопустимым в данном контексте.

\subsection{Синтаксические ошибки}

Пример:

\begin{lstlisting}
if (x > 0 {
    y = 1;
}
\end{lstlisting}

Пропущена закрывающая скобка.

\subsection{Семантические ошибки}

Пример:

\begin{lstlisting}
x = y + 1;
\end{lstlisting}

Если переменная \texttt{y} не объявлена, это семантическая ошибка.

Другой пример:

\begin{lstlisting}
int x;
x = "hello";
\end{lstlisting}

Здесь нарушена совместимость типов.

\section{Сравнение основных промежуточных представлений}

\[
\begin{array}{|p{4cm}|p{5cm}|p{5cm}|}
\hline
\textbf{Представление} & \textbf{Преимущества} & \textbf{Недостатки} \\
\hline
Последовательность символов
&
Прямо соответствует исходному файлу; удобно для чтения входа
&
Нет структуры программы
\\
\hline
Последовательность лексем
&
Убраны пробелы и комментарии; выделены значимые элементы
&
Нет синтаксического дерева и семантики
\\
\hline
Синтаксическое дерево
&
Полностью отражает грамматический вывод
&
Содержит много служебных узлов
\\
\hline
AST
&
Компактно отражает структуру программы
&
Не всегда удобно для глобальных оптимизаций
\\
\hline
Трёхадресный код
&
Удобен для оптимизаций и анализа потока данных
&
Менее близок к исходному языку
\\
\hline
SSA
&
Упрощает анализ определений и использований переменных
&
Требует $\varphi$-функций и преобразования обратно перед генерацией кода
\\
\hline
Низкоуровневое IR
&
Близко к машинному коду
&
Потеря высокоуровневой информации
\\
\hline
\end{array}
\]

\section{Итоговая схема уровней представления}

\[
\begin{array}{c}
\text{Символы}
\\
\downarrow
\\
\text{Лексемы}
\\
\downarrow
\\
\text{Синтаксическое дерево}
\\
\downarrow
\\
\text{AST}
\\
\downarrow
\\
\text{Высокоуровневое IR}
\\
\downarrow
\\
\text{Среднеуровневое IR}
\\
\downarrow
\\
\text{Низкоуровневое IR}
\\
\downarrow
\\
\text{Ассемблер}
\\
\downarrow
\\
\text{Машинный код}
\end{array}
\]

\section{Краткие выводы}

\begin{enumerate}
    \item Транслятор преобразует программу с исходного языка на целевой язык с сохранением семантики.
    \item Компилятор обычно состоит из фронтенда, мидлэнда и бэкенда.
    \item Фронтенд выполняет лексический, синтаксический и семантический анализ.
    \item Мидлэнд использует промежуточное представление для анализа и оптимизаций.
    \item Бэкенд генерирует код для конкретной целевой архитектуры.
    \item Промежуточные представления позволяют отделить исходный язык от целевой машины.
    \item На ранних этапах программа представляется как последовательность символов и лексем.
    \item Синтаксическое дерево отражает полный грамматический вывод, а AST --- только существенную структуру.
    \item Высокоуровневое IR близко к исходному языку, среднеуровневое удобно для оптимизаций, низкоуровневое близко к машинному коду.
    \item Важнейшие формы IR: деревья, трёхадресный код, квадрупли, трипли, стековый код, CFG и SSA.
    \item Оптимизации должны сохранять семантику программы.
    \item Многие оптимизации основаны на анализе потока данных и вычислении неподвижных точек.
\end{enumerate}

\end{document}
